% GARCIA SANGOQUIZA — CAP y PACELC; comparación pendiente.
\renewcommand{\fuentespie}{Fuentes: \cite{gsGilbertLynch2002,gsClabE1}}
\begin{frame}{CAP: coordinar datos al ejecutar en paralelo}
  \small
  \textbf{Con una partición de red no se garantizan C y A a la vez.}
  \begin{description}
    \item[C] Consistencia: operaciones como sobre una única copia actualizada.
    \item[A] Disponibilidad: cada petición a un nodo operativo se completa.
    \item[P] Tolerancia a la pérdida de comunicación entre nodos.
  \end{description}
  \begin{block}{Ejemplo general: dos réplicas, una última unidad}
    Aisladas, ambas aceptan una venta: se arriesga la sobreventa.
    Coordinar antes de confirmar puede impedir completar una venta.
  \end{block}
  \begin{block}{CLAB: dos nodos intentan reservar el mismo laboratorio}
    \textbf{Propuesta:} sin coordinación, el nodo aislado no confirma
    la reserva. Se prioriza C sobre A para esa operación.
  \end{block}
  {\scriptsize Escenarios ilustrativos; CLAB actual no implementa reservas
  ni réplicas. Fuentes: \cite{gsGilbertLynch2002,gsClabE1}.}
\end{frame}
\renewcommand{\fuentespie}{Fuentes: [PENDIENTE]}

\renewcommand{\fuentespie}{Fuentes: \cite{gsAbadi2012,gsMicrosoftConsistency2025}}
\begin{frame}{PACELC: coordinar también cuesta sin fallos}
  \small
  \begin{columns}[T,onlytextwidth]
    \begin{column}{0.48\textwidth}
      \begin{block}{Si hay partición (P)}
        Disponibilidad (A)\par
        frente a consistencia (C).
      \end{block}
    \end{column}
    \begin{column}{0.48\textwidth}
      \begin{block}{Si no la hay (E)}
        Menor latencia (L)\par
        frente a consistencia (C).
      \end{block}
    \end{column}
  \end{columns}
  \begin{block}{Ejemplo general: leer un precio replicado}
    Réplica cercana: menos espera posible, pero un valor antiguo.\par
    Lectura actualizada: puede exigir comunicación adicional.
  \end{block}
  \begin{block}{CLAB propuesto: consultas paralelas de disponibilidad}
    Consultar réplicas orienta; \textbf{no confirma una reserva}.\par
    La confirmación requiere validación y registro atómicos coordinados.
  \end{block}
  {\scriptsize Datos replicados hipotéticos; no implementados ni medidos en
  CLAB. Fuentes: \cite{gsAbadi2012,gsMicrosoftConsistency2025}.}
\end{frame}
\renewcommand{\fuentespie}{Fuentes: [PENDIENTE]}

\begin{frame}{Comparación de sistemas de cómputo paralelo}
  [PENDIENTE: contenido verificado y cita]\par
\end{frame}
